\documentclass[12pt]{article}
\usepackage[T1]{fontenc}
\usepackage[utf8]{inputenc}
\usepackage{lmodern}
\usepackage{textcomp}
\usepackage{enumitem}
\usepackage{geometry}
\geometry{margin=1in}
\begin{document}
\begin{center}
Answer Key - Chemistry - Undergraduate Level Assessment 21 \
\small (Answers are ordered exactly as in the question paper.)
\end{center}

\section*{Multiple Choice}
\begin{enumerate}[label=\arabic*.]
\item \textbf{Answer:} D
\\ \textit{Options:} A) BaH 2; B) BaOH; C) Ba; D) BaO
\\ \textit{Note:} The anhydride of Ba(OH)2 is BaO because it is formed by the removal of water (H$_{2}$O) from the hydroxide, resulting in the oxide of barium.
\item \textbf{Answer:} D
\\ \textit{Options:} A) B; B) Al; C) In; D) Tl
\\ \textit{Note:} The +1 oxidation state is more stable than the +3 oxidation state for thallium (Tl) due to the relativistic effects and the inert pair effect, which favor the stability of the fully filled 6 s subshell in Tl.
\item \textbf{Answer:} D
\\ \textit{Options:} A) racemic mixture; B) single pure enantiomer; C) mixture of two diastereomers in equal amounts; D) meso compound
\\ \textit{Note:} The reduction of D-xylose with NaBH 4 leads to the formation of a sugar alcohol that has two stereogenic centers and an internal plane of symmetry, classifying it as a meso compound.
\item \textbf{Answer:} B
\\ \textit{Options:} A) 0.0471 Hz; B) 21.2 Hz; C) 42.4 Hz; D) 66.7 Hz
\\ \textit{Note:} The linewidth can be calculated using the relationship Δν = 1/(2πT 2), and substituting T$_{2}$ = 15 ms yields a linewidth of approximately 21.2 Hz.
\item \textbf{Answer:} A
\\ \textit{Options:} A) hexane; B) 2-methylpentane; C) 3-methylpentane; D) 2,3-dimethylbutane
\\ \textit{Note:} Hexane has a linear structure that allows for three distinct types of carbon environments-primary, secondary, and tertiary-resulting in three distinct chemical shifts in its 13 C NMR spectrum.
\end{enumerate}

\section*{True / False}
\begin{enumerate}[label=\arabic*.]
\setcounter{enumi}{5}
\item \textbf{Answer:} True
\\ \textit{Options:} A) True; B) False
\\ \textit{Note:} The statement is true because the 55 Mn nuclide possesses three distinct nuclear energy levels due to its specific configuration of protons and neutrons, which allows for transitions between these states.
\item \textbf{Answer:} True
\\ \textit{Options:} A) True; B) False
\\ \textit{Note:} The equilibrium polarization of 13 C nuclei in a 20.0 T magnetic field at 300 K is true because the nuclear magnetic resonance polarization can be calculated using the Boltzmann distribution and is consistent with the measured value of 1.71 x 10^-5 at that temperature and field strength.
\item \textbf{Answer:} True
\\ \textit{Options:} A) True; B) False
\\ \textit{Note:} Arsenic-doped silicon is classified as an n-type semiconductor because arsenic introduces extra electrons into the silicon lattice, acting as a donor impurity that increases the material's electrical conductivity.
\item \textbf{Answer:} False
\\ \textit{Options:} A) True; B) False
\\ \textit{Note:} The infrared absorption frequency shift between DCl and HCl is primarily due to differences in their reduced masses rather than their dipole moments, as the mass difference affects the vibrational frequency of the molecules.
\item \textbf{Answer:} True
\\ \textit{Options:} A) True; B) False
\\ \textit{Note:} Silicon is abundant in the Earth's crust and primarily occurs as silicon dioxide (SiO 2) and in silicate minerals, which are the most prevalent forms of silicon in nature.
\end{enumerate}

\section*{Long Form Response}
\begin{enumerate}[label=\arabic*.]
\setcounter{enumi}{10}
\item \textbf{Answer:} Spin trapping is a preferred method for detecting free radical intermediates in chemical reactions primarily because it allows for the stabilization and identification of these transient species through the formation of spin adducts. This technique requires lower power than the direct detection of free radicals, making it more accessible and practical in many experimental setups. By using a spin trap, the free radicals react with the trap to form more stable adducts, which can then be detected using techniques like electron paramagnetic resonance (EPR) spectroscopy. This enhanced stability not only facilitates easier detection but also provides clearer insights into the reaction mechanisms involving free radicals. Overall, the advantages of spin trapping significantly improve the reliability and effectiveness of free radical detection in chemical research.
\\ \textit{Note:} correct because spin trapping stabilizes free radical intermediates by forming detectable spin adducts, making the detection process more practical and efficient compared to direct methods.
\item \textbf{Answer:} In a zero-order chemical reaction, the rate of the reaction is constant and is independent of the concentration of the reactants. The rate expression for a zero-order reaction is given by r = k, where r represents the rate of the reaction and k is the rate constant. This implies that the reaction proceeds at a steady rate until the reactants are depleted, regardless of their initial concentrations. The rate constant k, in this context, is a measure of how quickly the reaction occurs, and it remains constant over the course of the reaction as long as the temperature and other conditions are unchanged. Consequently, the implications for reaction kinetics indicate that zero-order reactions are often characterized by a linear decrease in reactant concentration over time, which can simplify the analysis of certain chemical processes.
\\ \textit{Note:} correct because it accurately describes that in a zero-order reaction, the reaction rate is constant and solely determined by the rate constant k, indicating that changes in reactant concentrations do not affect the rate until the reactants are exhausted.
\end{enumerate}

\vfill
\noindent\textit{End of Answer Key}
\end{document}